% =====================================================================
%  config/macros.tex — ký hiệu chuẩn của nhóm 10
%  Theo "bang_ky_hieu_chuan" (chốt 08/11/2026) và
%  "quy_uoc_chung_dang_he_thong" (chốt 06/11/2026):
%    - u, c, r, e, s là VECTƠ HÀNG; c = u G; s = r H^T;
%    - mọi phép cộng là cộng mod 2 (XOR), ký hiệu \oplus;
%    - dạng hệ thống: c = (u1 u2 u3 u4 p1 p2 p3);
%    - báo cáo đánh số bit 1..7, Python đánh số 0..6.
%  Dùng các lệnh dưới đây thay vì gõ tay để cả báo cáo thống nhất.
% =====================================================================

% Trường, không gian, phép toán
\newcommand{\GF}{\mathrm{GF}(2)}                 % GF(2)
\newcommand{\GFn}[1]{\mathrm{GF}(2)^{#1}}        % GF(2)^n
\newcommand{\xor}{\oplus}                         % cộng mod 2
\newcommand{\tp}{^{\mathsf{T}}}                   % chuyển vị: H\tp
\newcommand{\HT}{H^{\mathsf{T}}}                  % H^T
\newcommand{\PT}{P^{\mathsf{T}}}                  % P^T
\newcommand{\dmin}{d_{\min}}                      % khoảng cách tối thiểu
\newcommand{\wt}[1]{w(#1)}                        % trọng số Hamming
\newcommand{\dist}[2]{d(#1,#2)}                   % khoảng cách Hamming
\newcommand{\Hmr}{H'}                             % ma trận kiểm tra mở rộng 4x8
\DeclareMathOperator{\rank}{rank}
\newcommand{\vt}[1]{(#1)}                         % vectơ: \vt{1,0,1}
\newcommand{\bit}[1]{\texttt{#1}}                 % chuỗi bit trong văn bản
\newcommand{\code}[1]{\texttt{#1}}                % tên hàm, biến, file

% ---------------- Các ma trận đã chốt (dạng hệ thống) ---------------
% P (4x3)
\newcommand{\MatP}{\begin{bmatrix}
  1 & 1 & 0\\ 1 & 0 & 1\\ 0 & 1 & 1\\ 1 & 1 & 1
\end{bmatrix}}

% G = [I_4 | P] (4x7)
\newcommand{\MatG}{\left[\begin{array}{cccc|ccc}
  1 & 0 & 0 & 0 & 1 & 1 & 0\\
  0 & 1 & 0 & 0 & 1 & 0 & 1\\
  0 & 0 & 1 & 0 & 0 & 1 & 1\\
  0 & 0 & 0 & 1 & 1 & 1 & 1
\end{array}\right]}

% H = [P^T | I_3] (3x7)
\newcommand{\MatH}{\left[\begin{array}{cccc|ccc}
  1 & 1 & 0 & 1 & 1 & 0 & 0\\
  1 & 0 & 1 & 1 & 0 & 1 & 0\\
  0 & 1 & 1 & 1 & 0 & 0 & 1
\end{array}\right]}

% H^T (7x3): hàng thứ j là cột h_j của H
\newcommand{\MatHT}{\left[\begin{array}{ccc}
  1 & 1 & 0\\ 1 & 0 & 1\\ 0 & 1 & 1\\ 1 & 1 & 1\\ \hline
  1 & 0 & 0\\ 0 & 1 & 0\\ 0 & 0 & 1
\end{array}\right]}

% H' (4x8) của mã mở rộng (8,4)
\newcommand{\MatHmr}{\left[\begin{array}{ccccccc|c}
  1 & 1 & 0 & 1 & 1 & 0 & 0 & 0\\
  1 & 0 & 1 & 1 & 0 & 1 & 0 & 0\\
  0 & 1 & 1 & 1 & 0 & 0 & 1 & 0\\ \hline
  1 & 1 & 1 & 1 & 1 & 1 & 1 & 1
\end{array}\right]}

% H dạng cổ điển (cột j là số j viết nhị phân, hàng 1 là bit thấp nhất)
\newcommand{\MatHcd}{\begin{bmatrix}
  1 & 0 & 1 & 0 & 1 & 0 & 1\\
  0 & 1 & 1 & 0 & 0 & 1 & 1\\
  0 & 0 & 0 & 1 & 1 & 1 & 1
\end{bmatrix}}
